% =============================================================================
%  tesis.tex
%  -------------------------------------------------------------------------
%  Sesgo socioeconómico en modelos de clasificación automática de texto:
%  un enfoque sociotécnico para su detección y medición
%
%  Tesis de Licenciatura en Ciencias de la Computación
%  Autor:   Axel Casas Espinosa
%  Asesor:  M. en C.I.C. Manuel Soto Romero
% =============================================================================
%
%  COMPILACIÓN
%  -----------
%  Con bibtex (configuración por defecto, no requiere biber):
%      pdflatex tesis.tex
%      bibtex   tesis
%      pdflatex tesis.tex
%      pdflatex tesis.tex
%
%  Con latexmk (recomendado si está disponible):
%      latexmk -pdf tesis.tex
%
%  Con biber (recomendado, requiere editar preamble.tex):
%      pdflatex tesis.tex
%      biber    tesis
%      pdflatex tesis.tex
%      pdflatex tesis.tex
% =============================================================================

\documentclass[12pt,letterpaper,twoside,openright]{book}

% =============================================================================
%  preamble.tex
%  Configuración general de la tesis. Cargado desde tesis.tex.
% =============================================================================

% --- Codificación y tipografía ------------------------------------------------
\usepackage[utf8]{inputenc}
\usepackage[T1]{fontenc}
\usepackage{lmodern}                % Latin Modern (mejor que Computer Modern por defecto)
\usepackage{textcomp}
\usepackage{microtype}              % Mejoras tipográficas finas

% --- Idioma -------------------------------------------------------------------
\usepackage[spanish,es-tabla,es-nodecimaldot,es-noindentfirst]{babel}
\usepackage{csquotes}               % Comillas correctas; requerido por biblatex
\selectlanguage{spanish}

% --- Geometría de página (estándar de tesis) ---------------------------------
\usepackage[letterpaper,
            inner=3.5cm,            % margen interior (encuadernación)
            outer=2.5cm,
            top=3cm,
            bottom=3cm,
            headheight=15pt,
            headsep=1cm,
            footskip=1.2cm]{geometry}

% --- Interlineado e indentación ---------------------------------------------
\usepackage{setspace}
\onehalfspacing                     % 1.5 de interlineado (estándar de tesis)
\setlength{\parindent}{1.25em}
\setlength{\parskip}{0.2em plus 0.05em minus 0.05em}

% --- Matemáticas -------------------------------------------------------------
\usepackage{amsmath,amssymb,amsthm,mathtools}
\usepackage{bm}

% --- Gráficos, color y figuras ----------------------------------------------
\usepackage{graphicx}
\graphicspath{{imagenes/}}
\usepackage{xcolor}
\definecolor{azulUNAM}{HTML}{003B71}
\definecolor{oroUNAM}{HTML}{C8A047}
\definecolor{grisCodigo}{HTML}{F4F4F4}
\usepackage{caption}
\captionsetup{font=small,labelfont=bf,labelsep=period,format=plain}
\usepackage{subcaption}
\usepackage{float}

% --- Tablas ------------------------------------------------------------------
\usepackage{booktabs}               % \toprule, \midrule, \bottomrule
\usepackage{tabularx}
\usepackage{multirow}
\usepackage{array}

% --- Listas ------------------------------------------------------------------
\usepackage{enumitem}
\setlist{itemsep=0.25em,topsep=0.3em}

% --- Encabezados y pies ------------------------------------------------------
\usepackage{fancyhdr}
\usepackage{emptypage}              % Páginas en blanco verdaderamente en blanco
\pagestyle{fancy}
\fancyhf{}
\fancyhead[LE,RO]{\thepage}
\fancyhead[RE]{\nouppercase{\leftmark}}     % Capítulo
\fancyhead[LO]{\nouppercase{\rightmark}}    % Sección
\renewcommand{\headrulewidth}{0.4pt}
\renewcommand{\footrulewidth}{0pt}

\fancypagestyle{plain}{%
  \fancyhf{}%
  \fancyfoot[C]{\thepage}%
  \renewcommand{\headrulewidth}{0pt}%
}

% --- Formato de capítulos y secciones ---------------------------------------
\usepackage{titlesec}
\titleformat{\chapter}[display]
  {\normalfont\huge\bfseries\color{azulUNAM}}
  {\filright\Large\MakeUppercase{\chaptertitlename}\ \thechapter}
  {1ex}
  {\titlerule\vspace{1ex}\filright}
\titlespacing*{\chapter}{0pt}{-20pt}{30pt}

\titleformat{\section}
  {\normalfont\Large\bfseries\color{azulUNAM!90!black}}{\thesection}{0.75em}{}
\titleformat{\subsection}
  {\normalfont\large\bfseries}{\thesubsection}{0.6em}{}
\titleformat{\subsubsection}
  {\normalfont\normalsize\bfseries\itshape}{\thesubsubsection}{0.5em}{}

% --- Código fuente (listings) ------------------------------------------------
\usepackage{listings}
\lstdefinestyle{python}{
  language=Python,
  basicstyle=\ttfamily\footnotesize,
  keywordstyle=\color{azulUNAM}\bfseries,
  commentstyle=\color{gray!80!black}\itshape,
  stringstyle=\color{oroUNAM!80!black},
  backgroundcolor=\color{grisCodigo},
  numbers=left,
  numberstyle=\tiny\color{gray},
  numbersep=8pt,
  frame=single,
  rulecolor=\color{gray!40},
  framerule=0.4pt,
  showstringspaces=false,
  breaklines=true,
  breakatwhitespace=true,
  tabsize=4,
  captionpos=b,
  literate={á}{{\'a}}1 {é}{{\'e}}1 {í}{{\'i}}1 {ó}{{\'o}}1 {ú}{{\'u}}1
           {Á}{{\'A}}1 {É}{{\'E}}1 {Í}{{\'I}}1 {Ó}{{\'O}}1 {Ú}{{\'U}}1
           {ñ}{{\~n}}1 {Ñ}{{\~N}}1 {¿}{{?`}}1 {¡}{{!`}}1
}
\lstset{style=python}
\renewcommand{\lstlistingname}{Código}
\renewcommand{\lstlistlistingname}{Índice de códigos}

% --- Entornos de teoremas (en español) --------------------------------------
\theoremstyle{definition}
\newtheorem{definicion}{Definición}[chapter]
\newtheorem{ejemplo}{Ejemplo}[chapter]

\theoremstyle{plain}
\newtheorem{teorema}{Teorema}[chapter]
\newtheorem{proposicion}[teorema]{Proposición}
\newtheorem{lema}[teorema]{Lema}
\newtheorem{corolario}[teorema]{Corolario}

\theoremstyle{remark}
\newtheorem{observacion}{Observación}[chapter]

% --- Bibliografía (biblatex) -------------------------------------------------
% Si tienes biber instalado, cambia "backend=bibtex" por "backend=biber".
\usepackage[backend=bibtex,
            style=ieee,
            sorting=nyt,
            block=space,
            bibencoding=utf8]{biblatex}
\addbibresource{referencias.bib}

% --- Hipervínculos (debe cargarse casi al final) ----------------------------
\usepackage[hidelinks,
            colorlinks=false,
            pdfstartview=Fit,
            pdfusetitle,
            bookmarksnumbered=true]{hyperref}

% --- Comandos personalizados -------------------------------------------------
\newcommand{\INS}{\textsc{ins}\xspace}
\newcommand{\PLN}{\textsc{pln}\xspace}
\newcommand{\NLP}{\textsc{nlp}\xspace}
\usepackage{xspace}

% Cita en bloque destacada
\newenvironment{citablock}
  {\begin{quote}\itshape\small}
  {\end{quote}}


% --- Metadatos del PDF ------------------------------------------------------
\hypersetup{
  pdftitle  = {Sesgo socioeconómico en modelos de clasificación automática de texto},
  pdfauthor = {Axel Casas Espinosa},
  pdfsubject= {Tesis de Licenciatura en Ciencias de la Computación},
  pdfkeywords={sesgo algorítmico, procesamiento de lenguaje natural,
               discriminación lingüística, sociolectos, evaluación sociotécnica}
}

\begin{document}

% =============================================================================
%  Materia preliminar (numeración romana)
% =============================================================================
\frontmatter
\pagestyle{plain}

% =============================================================================
%  portada.tex
%  Portada en estilo Facultad de Ciencias / UNAM. Edita los campos según
%  tu institución y formato oficial.
% =============================================================================

\begin{titlepage}
\thispagestyle{empty}
\begin{center}

\vspace*{0.2cm}

{\Large\bfseries UNIVERSIDAD NACIONAL AUTÓNOMA \\[0.2em] DE MÉXICO}

\vspace{0.4cm}

{\large\bfseries FACULTAD DE CIENCIAS}

\vspace{1.0cm}

% Reemplaza con el logo institucional cuando lo tengas:
%   \includegraphics[width=4.5cm]{imagenes/logo_unam.png}
%   \hspace{1.2cm}
%   \includegraphics[width=3.8cm]{imagenes/logo_fc.png}
\rule{0pt}{4cm}\textit{[espacio reservado para escudos institucionales]}

\vspace{1.4cm}

{\Large\bfseries
Sesgo socioeconómico en modelos \\[0.25em]
de clasificación automática de texto: \\[0.25em]
un enfoque sociotécnico para \\[0.25em]
su detección y medición
}

\vspace{1.6cm}

{\LARGE\bfseries T\hspace{0.5em}E\hspace{0.5em}S\hspace{0.5em}I\hspace{0.5em}S}

\vspace{0.8cm}

QUE PARA OBTENER EL TÍTULO DE:

\vspace{0.4cm}

{\large\bfseries LICENCIADO EN CIENCIAS DE LA COMPUTACIÓN}

\vspace{1.0cm}

P\hspace{0.4em}R\hspace{0.4em}E\hspace{0.4em}S\hspace{0.4em}E\hspace{0.4em}N\hspace{0.4em}T\hspace{0.4em}A:

\vspace{0.4cm}

{\large\bfseries AXEL CASAS ESPINOSA}

\vspace{1.4cm}

\begin{minipage}{0.85\textwidth}
\centering
DIRECTOR DE TESIS: \\[0.4em]
{\large\bfseries M.\ en C.I.C.\ MANUEL SOTO ROMERO} \\[0.3em]
{\small Departamento de Matemáticas}
\end{minipage}

\vfill

CIUDAD UNIVERSITARIA, CD.\ MX., 2026

\end{center}
\end{titlepage}
\cleardoublepage


% =============================================================================
%  dedicatoria.tex
% =============================================================================

\thispagestyle{empty}
\vspace*{0.3\textheight}
\begin{flushright}
\itshape
[Espacio reservado para la dedicatoria.]
\end{flushright}
\cleardoublepage

% =============================================================================
%  agradecimientos.tex
% =============================================================================

\chapter*{Agradecimientos}
\addcontentsline{toc}{chapter}{Agradecimientos}
\markboth{Agradecimientos}{}

[Texto pendiente.]

\cleardoublepage

% =============================================================================
%  resumen.tex
% =============================================================================

\chapter*{Resumen}
\addcontentsline{toc}{chapter}{Resumen}
\markboth{Resumen}{}

Los sistemas de clasificación automática de texto basados en procesamiento de
lenguaje natural se emplean cada vez con mayor frecuencia en procesos de alto
impacto social, como el filtrado de currículums, la moderación de contenido o
la categorización de solicitudes en plataformas digitales. Aunque suelen
presentarse como herramientas neutrales, distintas investigaciones han mostrado
que pueden reproducir o amplificar desigualdades sociales preexistentes. Este
trabajo aborda, desde un enfoque sociotécnico, la pregunta de en qué medida
estos sistemas presentan sesgos asociados al \emph{registro lingüístico} o
\emph{sociolecto}, esto es, a variaciones de estilo que no afectan el contenido
sustantivo de la información evaluada.

La metodología combina revisión documental, recolección de percepciones
mediante encuesta a personas jóvenes de distintos contextos socioeconómicos en
México y experimentación computacional sobre uno o más modelos de clasificación
textual. A partir de los patrones identificados se construye un conjunto de
textos controlados que mantienen constante el contenido informativo y varían
únicamente el registro lingüístico, los cuales se utilizan como
\emph{benchmark}. Sobre los resultados experimentales se define y calcula un
indicador preliminar, el \emph{Índice de Neutralidad Sociolectal} (\INS),
orientado a medir el grado de consistencia con que un sistema responde a estas
variaciones.

\paragraph{Palabras clave.} sesgo algorítmico, procesamiento de lenguaje
natural, discriminación lingüística, sociolectos, evaluación sociotécnica,
equidad algorítmica.

\vspace{1.5cm}

\section*{Abstract}
\addcontentsline{toc}{section}{Abstract}

Automatic text classification systems based on natural language processing are
increasingly deployed in socially consequential pipelines, such as résumé
screening, content moderation, and request triage in digital platforms.
Although often framed as neutral tools, prior work has shown that they may
reproduce or amplify pre-existing social inequalities. This thesis addresses,
from a sociotechnical perspective, the extent to which such systems exhibit
biases associated with linguistic \emph{register} or \emph{sociolect}, that
is, stylistic variation that does not affect the substantive content of the
information being evaluated.

The methodology combines a documentary review, a perception survey administered
to young Mexicans from different socioeconomic backgrounds, and a computational
experiment over one or more text classification models. From the identified
patterns, a controlled benchmark is built in which informational content is
held constant while linguistic register is varied. Over the experimental
results, a preliminary indicator is defined and computed, the
\emph{Sociolectal Neutrality Index} (\INS), aimed at measuring the consistency
of system responses across such variations.

\paragraph{Keywords.} algorithmic bias, natural language processing, linguistic
discrimination, sociolects, sociotechnical evaluation, algorithmic fairness.

\cleardoublepage


\tableofcontents
\listoffigures
\listoftables

% =============================================================================
%  Cuerpo principal (numeración arábiga)
% =============================================================================
\mainmatter
\pagestyle{fancy}

% =============================================================================
%  Capítulo 1 -- Introducción
% =============================================================================
\chapter{Introducción}
\label{cap:introduccion}

\section{Contexto del problema}
\label{sec:contexto}

Los sistemas de clasificación automática de texto basados en técnicas de
procesamiento de lenguaje natural (\PLN) se utilizan cada vez con mayor
frecuencia en procesos como el filtrado automatizado de currículums, la
moderación de contenido o la categorización de información en plataformas
digitales. Aunque estos sistemas suelen presentarse como herramientas neutrales
orientadas a la eficiencia, diversas investigaciones han mostrado que pueden
reproducir o amplificar desigualdades sociales preexistentes cuando aprenden
de datos históricos o de patrones lingüísticos asociados a contextos sociales
específicos. En particular, los modelos de \PLN pueden incorporar sesgos
provenientes de los datos, del proceso de anotación o del diseño del sistema,
lo que puede generar resultados sistemáticamente desfavorables para ciertos
grupos sociales \cite{HovySpruit2016,HovyPrabhumoye2021}.

Desde la sociolingüística se ha documentado que el uso del lenguaje varía de
acuerdo con factores como el nivel educativo, la clase social o el contexto
cultural, y que estas variaciones pueden ser objeto de estigmatización o
discriminación. La \emph{discriminación lingüística} ocurre cuando determinados
registros, dialectos o estilos de habla se asocian con juicios negativos sobre
la capacidad, la inteligencia o la credibilidad de una persona, reproduciendo
jerarquías sociales preexistentes \cite{Drozdzowicz2024,Weissler2022}. En
entornos digitales mediados por algoritmos, estas asociaciones podrían
trasladarse a sistemas automatizados que procesan lenguaje natural, generando
respuestas diferenciadas ante variaciones sociolectales que no afectan el
contenido sustantivo de la información.

\section{Planteamiento del problema}
\label{sec:planteamiento}

A pesar del creciente uso de sistemas de \PLN en procesos con consecuencias
sociales relevantes, la mayoría de las evaluaciones de equidad disponibles en
la literatura se han centrado en atributos demográficos categóricos
---típicamente género, raza o edad--- y en lenguas distintas del español. La
dimensión \emph{sociolectal}, entendida aquí como variación lingüística
asociada a factores socioeconómicos, ha recibido menos atención sistemática,
particularmente en el español hablado en México. Como consecuencia, no se
dispone de instrumentos estandarizados para evaluar si un clasificador de
texto responde de manera consistente cuando se le presenta el mismo contenido
informativo expresado en registros distintos.

Este vacío instrumental tiene implicaciones prácticas. En procesos como el
filtrado automatizado de currículums, un comportamiento sistemáticamente
distinto del modelo ante variaciones de registro puede traducirse en
desventajas concretas para personas cuyo estilo de escritura difiere del
registro dominante en los datos de entrenamiento, sin que ello refleje
diferencias reales en el contenido evaluado.

\section{Pregunta de investigación}
\label{sec:pregunta}

\begin{citablock}
¿En qué medida los sistemas de clasificación automatizada de texto presentan
sesgos asociados al registro lingüístico o sociolectal, y cómo puede diseñarse
un indicador que permita evaluar su neutralidad frente a estas variaciones?
\end{citablock}

\section{Objetivos}
\label{sec:objetivos}

\subsection*{Objetivo general}

Analizar la presencia de sesgos asociados al registro lingüístico o sociolectal
en sistemas de clasificación automatizada de texto, mediante la construcción de
textos controlados que varían únicamente en su estilo lingüístico, con el
propósito de evaluar su comportamiento frente a estas variaciones y proponer un
indicador denominado \emph{Índice de Neutralidad Sociolectal} (\INS) para medir
el grado de equidad lingüística en dichos sistemas.

\subsection*{Objetivos específicos}

\begin{enumerate}[label=\textbf{O\arabic*.},leftmargin=2.2em]
  \item Documentar, mediante revisión de literatura académica y reportes
    institucionales, cómo el clasismo estructural y la discriminación
    lingüística se manifiestan en el contexto mexicano y de qué manera podrían
    reflejarse en sistemas automatizados de procesamiento de lenguaje.
  \item Diseñar un conjunto de textos o currículums simulados en los que el
    contenido informativo permanezca constante, variando únicamente el
    registro lingüístico, con el fin de construir un conjunto de pruebas
    controladas para evaluar sistemas de clasificación automática de texto.
  \item Evaluar el comportamiento de uno o más modelos de clasificación
    textual frente a las variaciones lingüísticas definidas, comparando sus
    resultados para identificar posibles diferencias sistemáticas asociadas
    al estilo de escritura.
  \item Diseñar y calcular un indicador denominado \emph{Índice de Neutralidad
    Sociolectal} (\INS) que permita medir el grado de consistencia o variación
    en las respuestas del sistema ante cambios en el registro lingüístico.
\end{enumerate}

\section{Alcances y limitaciones}
\label{sec:alcances}

\paragraph{Alcances.} El estudio se restringe a sistemas de clasificación
automática de texto sobre lengua española, con énfasis en el español hablado
en México. La evaluación experimental se circunscribe a uno o más modelos de
clasificación seleccionados explícitamente en el Capítulo~\ref{cap:metodologia}
y al \emph{benchmark} controlado descrito en el Capítulo~\ref{cap:resultados}.

\paragraph{Limitaciones.} Las limitaciones principales del trabajo son tres.
Primero, el alcance muestral de la encuesta de percepciones no permite
inferencias representativas a nivel poblacional; sus resultados se interpretan
como insumo cualitativo para el diseño del \emph{benchmark}. Segundo, el
\emph{benchmark} controlado opera sobre textos simulados y no captura la
totalidad de la variación sociolectal observada en datos reales. Tercero, el
\INS se propone como indicador \emph{preliminar} y su validación empírica
exhaustiva se identifica como trabajo futuro.

\section{Contribuciones de la tesis}
\label{sec:contribuciones}

Este trabajo aporta los siguientes elementos.

\begin{itemize}
  \item Una articulación entre la literatura sobre discriminación lingüística
    y la de evaluación de sesgos en \PLN, situada en el contexto mexicano.
  \item Un \emph{benchmark} controlado de textos en distintos registros
    sociolectales, con contenido informativo equivalente.
  \item Un protocolo experimental reproducible para evaluar la consistencia
    de un clasificador frente a variaciones de registro.
  \item La definición operativa y el cálculo de un indicador, el \INS, sobre
    los resultados experimentales.
\end{itemize}

\section{Estructura del documento}
\label{sec:estructura}

El resto de la tesis se organiza como sigue. El Capítulo~\ref{cap:marco}
presenta el marco teórico, articulando los conceptos de sesgo algorítmico,
clasificación automática de texto, discriminación lingüística y clasismo
estructural en el contexto mexicano. El Capítulo~\ref{cap:metodologia}
detalla la metodología sociotécnica empleada. El Capítulo~\ref{cap:resultados}
reporta los resultados de la encuesta de percepciones y de los experimentos
computacionales. El Capítulo~\ref{cap:ins} define formalmente el \INS,
describe su procedimiento de cálculo y discute su interpretación. El
Capítulo~\ref{cap:discusion} discute las implicaciones sociotécnicas de los
resultados y las limitaciones metodológicas del estudio. El
Capítulo~\ref{cap:conclusiones} presenta las conclusiones generales y
delinea líneas de investigación futura.

% =============================================================================
%  Capítulo 2 -- Marco teórico
%  -------------------------------------------------------------------------
%  NOTA PARA AXEL:
%  - Este capítulo está redactado con contenido real, pero deliberadamente
%    sintético: cada sección es un primer borrador sólido que debes expandir
%    con detalle de las fuentes primarias.
%  - Las marcas «[VERIFICAR: ...]» señalan afirmaciones cuya formulación
%    exacta conviene cotejar contra el artículo citado antes de la entrega.
%  - Las marcas «[FUENTE MX: ...]» señalan dónde DEBES añadir bibliografía
%    mexicana específica (esto cubre tu objetivo específico O1). No inventé
%    referencias para México: no existen aún en referencias.bib.
% =============================================================================
\chapter{Marco teórico}
\label{cap:marco}

Este capítulo construye el andamiaje conceptual de la tesis articulando cuatro
cuerpos de literatura que rara vez se leen en conjunto. La
Sección~\ref{sec:sesgo-algoritmico} sitúa la noción de sesgo en sistemas de
inteligencia artificial y precisa el sentido en que se usará a lo largo del
trabajo. La Sección~\ref{sec:pln} introduce los elementos técnicos de la
clasificación automática de texto necesarios para interpretar los experimentos
posteriores. La Sección~\ref{sec:sociolectos} recupera, desde la
sociolingüística, los conceptos de variación sociolectal y discriminación
lingüística. Finalmente, la Sección~\ref{sec:clasismo} aterriza estas
discusiones en el contexto mexicano, que es el escenario empírico de la
investigación.

\section{Sesgo algorítmico en sistemas de inteligencia artificial}
\label{sec:sesgo-algoritmico}

\subsection{Tres sentidos de ``sesgo''}
\label{sec:tres-sentidos}

El término \emph{sesgo} es notoriamente polisémico, y buena parte de la
confusión en los debates públicos sobre inteligencia artificial proviene de no
distinguir sus acepciones. Conviene separar al menos tres.

En primer lugar, el \emph{sesgo estadístico} designa la diferencia sistemática
entre el valor esperado de un estimador y el parámetro que pretende estimar. Es
una noción técnica, neutra desde el punto de vista moral, y constitutiva de
todo aprendizaje: el llamado \emph{compromiso sesgo--varianza} muestra que
cierto sesgo es incluso deseable para controlar la varianza de un modelo. En
segundo lugar, el \emph{sesgo cognitivo} se refiere a las desviaciones
sistemáticas del juicio humano respecto de la racionalidad normativa, estudiadas
por la psicología. En tercer lugar, el \emph{sesgo social} o
\emph{discriminatorio} designa el trato sistemáticamente desigual de grupos
sociales por parte de un sistema, en formas que reproducen o amplifican
desigualdades preexistentes \cite{Mehrabi2021Survey}.

Esta tesis se ocupa exclusivamente del tercer sentido, cuya relevancia para el
\PLN fue señalada tempranamente por Hovy y Spruit \cite{HovySpruit2016} al
argumentar que las tecnologías del lenguaje tienen un impacto social propio
---sobre la exclusión, la sobregeneralización y la exposición de grupos---
que la comunidad técnica había tendido a ignorar. La distinción entre los tres
sentidos no es ociosa: un sistema puede estar estadísticamente bien calibrado
y, sin embargo, ser socialmente sesgado, porque los datos sobre los que aprende
ya codifican desigualdades del mundo. La cuestión, por tanto, no es meramente técnica sino
\emph{sociotécnica}: requiere examinar conjuntamente el comportamiento del
modelo y el contexto social en el que se inserta.

\subsection{Definición operativa}
\label{sec:definicion-sesgo}

Para que el problema sea evaluable empíricamente se adopta la siguiente
caracterización operativa.

\begin{definicion}[Sesgo social en un sistema de clasificación]
\label{def:sesgo}
Sea $f : \mathcal{X} \to \mathcal{Y}$ un clasificador y sea $A$ un atributo
protegido o socialmente sensible definido sobre la población de entradas. Se
dice que $f$ exhibe \emph{sesgo social} respecto a $A$ si existen valores
$a, a' \in A$ y una métrica de comportamiento $\mu$ tales que
\[
  \mu\bigl(f \mid A = a\bigr) \;\neq\; \mu\bigl(f \mid A = a'\bigr),
\]
y dicha diferencia no se justifica por una variación correspondiente en la
variable objetivo $\mathcal{Y}$ atribuible legítimamente a $A$.
\end{definicion}

La cláusula final ---``no se justifica por una variación legítima''--- es la que
carga con todo el peso normativo y, a la vez, la más difícil de operacionalizar.
En el caso de esta tesis, el atributo $A$ es el \emph{registro lingüístico} de
un texto, y la justificación se vuelve especialmente tenue: si dos textos
comunican el mismo contenido sustantivo y difieren solo en su estilo, no existe
ninguna variación legítima en $\mathcal{Y}$ que pueda explicar un trato
diferenciado. Es precisamente este caso límite el que el \emph{Índice de
Neutralidad Sociolectal} (\INS), definido en el Capítulo~\ref{cap:ins}, busca
capturar.

\subsection{Fuentes de sesgo en sistemas de procesamiento de lenguaje natural}
\label{sec:fuentes-sesgo}

Hovy y Prabhumoye \cite{HovyPrabhumoye2021} ofrecen una taxonomía influyente que
identifica cinco puntos del ciclo de vida de un sistema de \PLN en los que puede
introducirse sesgo social. Conviene revisarlos porque ofrecen un mapa de dónde
podría originarse el sesgo sociolectal que la tesis investiga.

\subsubsection{Sesgo en los datos}

Los modelos de \PLN aprenden de corpus que son muestras del lenguaje producido
en condiciones sociales concretas, no espejos neutrales de la lengua. El sesgo
de los datos surge cuando ciertos grupos, registros o variedades están
sobre- o subrepresentados, o cuando el texto disponible refleja desigualdades
históricas. A medida que los corpus de entrenamiento crecen hasta volúmenes que
ningún equipo puede auditar exhaustivamente, este problema se agrava: Bender
y colaboradores \cite{Bender2021Parrots} advierten que el tamaño no garantiza
diversidad y que los corpus masivos tienden a sobrerrepresentar las variedades
hegemónicas y los puntos de vista dominantes. [VERIFICAR: matiz exacto del
argumento sobre tamaño vs.\ diversidad.]

\subsubsection{Sesgo en las anotaciones}

Gran parte de los sistemas de clasificación se entrena con etiquetas asignadas
por anotadores humanos, cuyas intuiciones no son neutrales. La composición
demográfica del panel de anotación, las guías de etiquetado y los juicios
implícitos sobre qué cuenta como ``correcto'' o ``ofensivo'' se transfieren al
modelo. Un ejemplo documentado y directamente pertinente para esta tesis es el
de Sap y colaboradores \cite{Sap2019Risk}, quienes mostraron que los anotadores
tendían a clasificar como ofensivo el inglés afroamericano vernáculo con mayor
frecuencia que el inglés estándar, sesgo que los clasificadores entrenados sobre
esas anotaciones aprendían y reproducían. El caso ilustra el mecanismo central
del presente trabajo: una variedad lingüística asociada a un grupo social
recibe un trato sistemáticamente distinto que no responde al contenido del
mensaje.

\subsubsection{Sesgo en las representaciones lingüísticas}

Antes de ser clasificado, el texto se transforma en representaciones numéricas.
Estas representaciones no son inertes: codifican regularidades estadísticas del
corpus, incluidas las asociaciones estereotipadas. Caliskan y colaboradores
\cite{Caliskan2017Semantics} demostraron que los \emph{embeddings} de palabras
entrenados sobre texto de la web reproducen sesgos humanos medibles, análogos a
los que captan las pruebas de asociación implícita en psicología. Bolukbasi y
colaboradores \cite{Bolukbasi2016Man} ilustraron el fenómeno con analogías de
género estereotipadas y propusieron técnicas tempranas de mitigación. Estos
trabajos establecen que el sesgo puede residir en la capa de representación,
con independencia del clasificador final.

\subsubsection{Sesgo en el diseño del modelo}

La arquitectura, la función objetivo y los hiperparámetros de un modelo también
introducen sesgo, en particular porque los modelos tienden a \emph{amplificar}
las regularidades del entrenamiento más que a reproducirlas fielmente. Un
patrón presente en una minoría de los datos puede volverse la respuesta por
defecto del modelo si reduce la pérdida agregada. Para esta tesis, la
implicación es metodológica: el sesgo sociolectal observado en un clasificador
no puede atribuirse sin más a sus datos de entrenamiento, pues la magnitud del
efecto depende también de la familia de modelo; de ahí que el diseño
experimental del Capítulo~\ref{cap:metodologia} contemple modelos de
complejidad creciente.

\subsubsection{Sesgo en el diseño de la investigación}

La fuente más sutil es la elección misma de qué problemas se estudian, sobre qué
lenguas y para qué poblaciones. La concentración del \PLN en el inglés y en
contextos del Norte global determina qué fenómenos se vuelven visibles y cuáles
permanecen sin estudiar. Esta tesis es, en parte, una respuesta a ese sesgo de
agenda: investiga la variación sociolectal del español mexicano, un objeto
sistemáticamente ausente de los \emph{benchmarks} de equidad disponibles.

\subsection{Del sesgo a la equidad: una nota crítica}
\label{sec:critica-bias}

Blodgett y colaboradores \cite{Blodgett2020Survey} señalan que buena parte de la
literatura sobre ``sesgo'' en \PLN adolece de imprecisión conceptual: se invoca
el término sin especificar qué comportamiento es perjudicial, para quién y por
qué. Recomiendan anclar todo estudio de sesgo en un daño concreto y en su
relación con las desigualdades sociales que lo motivan. Esta tesis adopta esa
recomendación: el daño concreto es el trato diferenciado de hablantes según su
registro en procesos de consecuencias materiales ---como el filtrado de
currículums--- y la desigualdad de fondo es el clasismo lingüístico discutido
en la Sección~\ref{sec:clasismo}.

\section{Procesamiento de lenguaje natural y clasificación automática de texto}
\label{sec:pln}

\subsection{La tarea de clasificación de texto}
\label{sec:tarea-clasificacion}

La clasificación de texto es la tarea de asignar a un documento $x \in
\mathcal{X}$ una etiqueta $y \in \mathcal{Y}$ de un conjunto finito y
predefinido. Formalmente, se busca aprender una función $f : \mathcal{X} \to
\mathcal{Y}$ a partir de un conjunto de ejemplos etiquetados. La canalización
(\emph{pipeline}) típica comprende cuatro etapas: (i)~\emph{preprocesamiento}
del texto crudo; (ii)~\emph{representación} del documento como un vector o una
secuencia de vectores; (iii)~\emph{modelado}, donde un algoritmo de aprendizaje
estima una función de decisión; y (iv)~\emph{decisión}, donde se emite la
etiqueta y, opcionalmente, una distribución de probabilidad sobre las clases.
El \INS, definido más adelante, opera sobre las salidas de las etapas (iii) y
(iv): la decisión y la probabilidad posterior.

\subsection{Representaciones del texto}
\label{sec:representaciones}

La forma de representar el texto determina qué información queda disponible para
el clasificador. Tres familias resultan relevantes para esta tesis.

Las representaciones \emph{de bolsa de palabras} y su ponderación
\emph{TF--IDF} describen un documento por las frecuencias de los términos que
contiene, ignorando el orden. Son interpretables y computacionalmente baratas,
pero ciegas al contexto: dos textos con el mismo contenido y distinto registro
pueden recibir representaciones muy distintas solo por su léxico, lo que las
hace un punto de partida útil para estudiar el sesgo sociolectal.

Los \emph{embeddings} estáticos asignan a cada palabra un vector fijo aprendido
de grandes corpus, de modo que palabras de uso similar quedan próximas en el
espacio vectorial. Como se discutió en la Sección~\ref{sec:fuentes-sesgo}, estos
vectores codifican asociaciones sociales del corpus de origen.

Los \emph{embeddings} contextuales, popularizados por arquitecturas de tipo
\emph{transformer} como BERT \cite{Devlin2019BERT}, asignan a cada aparición de
una palabra un vector que depende de su contexto, lo que mejora notablemente el
desempeño en la mayoría de las tareas. Para el español existen modelos
preentrenados específicos, como BETO \cite{CanetePilar2020BETO}, que serán
candidatos naturales para la evaluación experimental del
Capítulo~\ref{cap:resultados}.

\subsection{Familias de modelos relevantes}
\label{sec:familias-modelos}

Para los fines de la tesis conviene distinguir tres familias de clasificadores,
con grados crecientes de complejidad y opacidad. Los \emph{modelos lineales}
(regresión logística, máquinas de vectores de soporte) sobre representaciones
TF--IDF ofrecen interpretabilidad y un comportamiento predecible. Los
\emph{modelos basados en \emph{transformers} preentrenados} y ajustados a la
tarea, como BETO \cite{CanetePilar2020BETO}, alcanzan mayor desempeño a costa de
transparencia. Por último, los \emph{modelos de lenguaje grandes} empleados como
clasificadores en modo \emph{zero-shot} o \emph{few-shot} permiten clasificar sin
entrenamiento específico, pero su comportamiento ante variaciones de registro
es poco estudiado y especialmente difícil de auditar \cite{Bender2021Parrots}.
La inclusión de al menos un representante de cada familia permitirá observar si
el sesgo sociolectal varía con la complejidad del modelo.

\subsection{Métricas estándar de evaluación y sus límites}
\label{sec:metricas}

El desempeño de un clasificador suele resumirse con la exactitud, la precisión,
la exhaustividad (\emph{recall}) y su media armónica, la medida F1, además de
métricas de calibración cuando interesan las probabilidades. Estas métricas son
\emph{agregadas}: describen el comportamiento promedio sobre un conjunto de
prueba. Su limitación, central para esta tesis, es que un modelo puede exhibir
una F1 alta y, simultáneamente, comportarse de manera sistemáticamente distinta
ante subgrupos o registros, porque los errores se compensan en el agregado
\cite{Blodgett2020Survey, Mehrabi2021Survey}. De ahí la necesidad de métricas
\emph{desagregadas} y de indicadores específicos de consistencia como el \INS,
que no miden cuánto acierta el sistema sino cuán establemente se comporta ante
variaciones que no deberían importar.

\section{Discriminación lingüística y sociolectos}
\label{sec:sociolectos}

\subsection{Variación sociolectal: conceptos básicos}
\label{sec:variacion}

La sociolingüística ha documentado de manera consistente que el lenguaje varía
de forma sistemática según factores sociales. Conviene precisar tres términos
que se usarán con cuidado en lo que sigue. Un \emph{dialecto} es una variedad
asociada principalmente a un grupo geográfico. Un \emph{sociolecto} es una
variedad asociada a un grupo social ---definido por clase, nivel educativo,
ocupación u otros factores estructurales. Un \emph{registro} es la variación que
un mismo hablante despliega según la situación comunicativa, desde lo más formal
a lo más coloquial. Estas dimensiones se entrecruzan: un rasgo puede ser a la
vez marcador sociolectal y de registro.

Lo crucial para esta tesis es que la variación sociolectal es, en lo esencial,
variación de \emph{forma} y no de \emph{contenido}: dos hablantes de
sociolectos distintos pueden comunicar exactamente la misma información. Cuando
un sistema responde de manera diferente a ambos, la diferencia no puede
atribuirse al mensaje, sino al cómo se dice ---es decir, a un marcador social.

\subsection{La discriminación lingüística como fenómeno social}
\label{sec:discriminacion-ling}

La \emph{discriminación lingüística} ocurre cuando determinados registros,
dialectos o sociolectos se asocian con juicios negativos sobre la capacidad, la
inteligencia o la credibilidad de quien los emplea, de modo que la forma del
habla se convierte en vía para reproducir jerarquías sociales preexistentes
\cite{Weissler2022}. Drożdżowicz \cite{Drozdzowicz2024} subraya que el fenómeno
es más complejo de lo que sugiere su nombre, pues entrelaza juicios estéticos,
epistémicos y morales que resultan difíciles de separar analíticamente.
[VERIFICAR: caracterización precisa de la tesis de Drożdżowicz; cotejar contra
el artículo antes de la entrega.] Lo que importa retener es que la
discriminación lingüística no es una reacción a deficiencias comunicativas
reales ---el mensaje se entiende perfectamente--- sino una sanción social
adherida a marcadores de grupo.

[FUENTE MX: aquí conviene introducir el concepto de \emph{ideología de la
lengua estándar} y vincularlo con la tradición sociolingüística clásica sobre
prestigio y estigma. Añade referencias canónicas de sociolingüística variacionista
y, si es posible, trabajos sobre el español.]

\subsection{Del juicio humano al sesgo algorítmico}
\label{sec:puente}

El argumento que articula esta tesis puede enunciarse como un puente en tres
pasos. Primero, existe una discriminación lingüística humana documentada que
sanciona ciertos sociolectos (Sección~\ref{sec:discriminacion-ling}). Segundo,
los sistemas de \PLN aprenden de texto producido y anotado por humanos, de modo
que pueden internalizar esas sanciones (Sección~\ref{sec:fuentes-sesgo}).
Tercero, existe evidencia empírica de que esto efectivamente ocurre: el ya
mencionado estudio de Sap y colaboradores \cite{Sap2019Risk} muestra un caso en
el que una variedad lingüística asociada a un grupo social recibe un trato
algorítmico sistemáticamente desfavorable. La hipótesis de trabajo de la tesis
es que un mecanismo análogo podría operar para los sociolectos del español
mexicano, hipótesis que el diseño experimental de los
Capítulos~\ref{cap:metodologia} y~\ref{cap:resultados} pone a prueba de forma
controlada.

\section{Clasismo estructural y lenguaje en el contexto mexicano}
\label{sec:clasismo}

\subsection{Clasismo estructural en México}
\label{sec:clasismo-mx}

El clasismo no se reduce a actitudes individuales de desprecio: es
\emph{estructural} en la medida en que se inscribe en instituciones, prácticas y
representaciones que distribuyen de manera desigual el reconocimiento y las
oportunidades según la posición social percibida. En el caso mexicano, el
clasismo se entrecruza además con marcadores raciales y geográficos, de modo que
clase, color de piel y origen tienden a leerse conjuntamente.

[FUENTE MX: desarrolla esta subsección con literatura mexicana específica
---estudios sobre discriminación (p.\ ej., informes de instituciones nacionales
dedicadas a prevenir la discriminación), trabajos sociológicos y antropológicos
sobre clasismo y colorismo. Esta subsección es el núcleo de tu objetivo
específico O1 y requiere fuentes que aún no están en referencias.bib. No incluí
citas aquí para no inventarlas.]

\subsection{Marcadores sociolingüísticos del clasismo en el español mexicano}
\label{sec:marcadores-mx}

En el español mexicano, ciertos rasgos lingüísticos funcionan como índices de
clase y son objeto de juicio social. Entre ellos suelen señalarse el contraste
entre léxico formal y coloquial, el uso de anglicismos como marca de cierto
capital cultural, fenómenos de pronunciación, marcadores discursivos
característicos y, en el registro escrito informal, la variación ortográfica.
Los términos \emph{naco} y \emph{fresa}, ampliamente usados en el habla
cotidiana, condensan precisamente la asociación entre formas de hablar (y de
presentarse) y posición social percibida, operando como etiquetas
estigmatizantes y de prestigio respectivamente.

Para los fines de esta tesis, la importancia de estos marcadores es metodológica:
constituyen las dimensiones de variación que el \emph{benchmark} controlado del
Capítulo~\ref{cap:metodologia} deberá operacionalizar, manteniendo constante el
contenido informativo. La selección final de los marcadores se justificará
empíricamente a partir de la encuesta de percepciones.

[FUENTE MX: respalda la caracterización de estos marcadores con trabajos
lingüísticos sobre el español de México; evita apoyarte solo en el conocimiento
común.]

\subsection{Implicaciones para sistemas automatizados desplegados en México}
\label{sec:implicaciones-mx}

Si los marcadores anteriores estructuran el juicio social humano y los sistemas
de \PLN aprenden de texto humano, cabe esperar que un clasificador entrenado
o ajustado con datos del español mexicano herede, al menos en parte, esas
asociaciones. El riesgo es concreto: en procesos como el filtrado automatizado
de currículums, una respuesta sistemáticamente distinta del modelo ante el
mismo contenido expresado en un registro estigmatizado se traduciría en
desventajas materiales para personas cuyo estilo de escritura se aparta del
registro dominante. Documentar y medir ese riesgo es el propósito de los
capítulos siguientes.

\section*{Síntesis del capítulo}
\addcontentsline{toc}{section}{Síntesis del capítulo}

El recorrido de este capítulo puede condensarse en cuatro proposiciones que
sostienen el resto de la tesis. Primera: el sesgo que interesa aquí es el
social, no el estadístico, y se vuelve evaluable cuando se fija un atributo
sensible y una métrica de comportamiento (Definición~\ref{def:sesgo}). Segunda:
ese sesgo puede entrar a un sistema de \PLN por cinco vías documentadas
---datos, anotaciones, representaciones, modelo y agenda de investigación---,
todas plausibles para el caso sociolectal. Tercera: la variación sociolectal es
variación de forma y no de contenido, de modo que un trato algorítmico
diferenciado ante ella no admite justificación informativa. Cuarta: en el
contexto mexicano existen marcadores lingüísticos de clase socialmente
sancionados, lo que hace verosímil ---aunque no demostrado--- que los sistemas
desplegados en ese contexto los hereden. Establecer si esto ocurre, y con qué
magnitud, es una pregunta empírica; el siguiente capítulo describe el diseño
con el que se aborda.

% =============================================================================
%  Capítulo 3 -- Metodología
% =============================================================================
\chapter{Metodología}
\label{cap:metodologia}

\section{Diseño general de la investigación}
\label{sec:diseno-general}

La investigación adopta una metodología de \emph{evaluación sociotécnica} que
combina tres componentes complementarios: (i) un análisis documental, (ii) la
recolección de percepciones sociales mediante encuesta y (iii) un experimento
computacional controlado. El supuesto que articula los tres componentes es
que la evaluación de sistemas algorítmicos desplegados en contextos sociales
concretos no puede limitarse a métricas técnicas abstractas, sino que requiere
contrastar el comportamiento del sistema con el fenómeno social que podría
estar replicando.

\begin{figure}[H]
  \centering
  % Reemplaza este placeholder con un diagrama del flujo metodológico.
  \fbox{\parbox{0.9\textwidth}{\centering\vspace{1cm}
    \textit{[Diagrama del flujo metodológico:}\\
    \textit{revisión documental} $\longrightarrow$ \textit{encuesta de
    percepciones} $\longrightarrow$ \textit{diseño del benchmark} $\longrightarrow$
    \textit{experimentos} $\longrightarrow$ \textit{cálculo del \INS\textit{]}}
    \vspace{1cm}}}
  \caption[Esquema general de la metodología sociotécnica.]{Esquema general
  de la metodología sociotécnica empleada.}
  \label{fig:flujo-metodologico}
\end{figure}

\section{Revisión documental}
\label{sec:revision-documental}

\subsection{Criterios de inclusión y exclusión}

[Bases bibliográficas consultadas; años; idiomas; tipos de documento;
criterios para incluir y excluir referencias.]

\subsection{Estrategia de búsqueda}

[Cadenas de búsqueda, palabras clave en español e inglés, número de resultados
recuperados, filtrado y selección final.]

\subsection{Síntesis temática}

[Descripción del procedimiento para sintetizar la literatura encontrada en
los ejes temáticos del Capítulo~\ref{cap:marco}.]

\section{Recolección de percepciones sobre discriminación tecnológica}
\label{sec:encuesta}

El segundo componente de la metodología consistió en una encuesta dirigida a
personas jóvenes residentes en México, orientada a recoger percepciones y
experiencias sobre el trato diferenciado ---humano y automatizado--- asociado
a la forma de hablar o escribir. La encuesta cumple un doble propósito dentro
del diseño sociotécnico. Por una parte, aporta evidencia situada sobre cómo se
experimenta la discriminación lingüística en el contexto mexicano,
complementando la revisión documental del objetivo~O1. Por otra, sus
respuestas abiertas constituyen el insumo empírico del que se deriva la
tipología de variaciones lingüísticas que estructura el \emph{benchmark}
(Sección~\ref{sec:benchmark-diseno}): en lugar de postular desde el escritorio
qué rasgos del registro escrito son objeto de juicio social, la tipología se
ancla en los rasgos que los propios participantes reportan como sancionados.

\subsection{Población y muestreo}
\label{sec:poblacion-muestreo}

La población objetivo se definió como personas de 18 a 30 años residentes en
México [DECIDIR: confirmar rango de edad con el asesor]. La elección de este
rango respondió a dos criterios. Primero, se trata del grupo con mayor
exposición a procesos automatizados de evaluación de texto ---postulaciones de
empleo y becas en línea, plataformas digitales, sistemas de atención
automatizada---, por lo que sus percepciones refieren a experiencias recientes
y no hipotéticas. Segundo, restringir la población a personas mayores de edad
evita el tratamiento de datos de menores y los requisitos de consentimiento
parental asociados.

El muestreo fue no probabilístico, por conveniencia y con propagación por bola
de nieve: el cuestionario se difundió a través de los contactos del
investigador y se pidió a cada participante compartirlo con otras personas del
rango de edad. Para mitigar el sesgo de red ---la tendencia de este tipo de
muestreo a concentrarse en el círculo inmediato del investigador, en este caso
universitario---, la difusión inicial se dirigió deliberadamente a contactos
de distintos niveles educativos y contextos socioeconómicos
[PENDIENTE: describir brevemente los canales de difusión efectivamente
utilizados].

Esta estrategia de muestreo impone una limitación que conviene dejar asentada
desde ahora: la muestra resultante no es representativa de la población joven
mexicana, por lo que los resultados cuantitativos de la encuesta no admiten
inferencia poblacional. Su función en esta investigación es de carácter
exploratorio e instrumental ---identificar dimensiones relevantes del fenómeno
y alimentar el diseño del \emph{benchmark}---, no estimar prevalencias. Esta
limitación se retoma en la Sección~\ref{sec:limitaciones}.
[PENDIENTE: reportar el tamaño muestral final alcanzado, $n = \ldots$, y el
periodo de levantamiento.]

\subsection{Instrumento}
\label{sec:instrumento}

El instrumento fue un cuestionario en línea, autoadministrado y anónimo, con
una duración estimada de respuesta de 10 a 12 minutos, estructurado en cuatro
secciones: consentimiento informado, datos sociodemográficos, percepciones en
escala tipo Likert y preguntas abiertas. El instrumento completo se reproduce
en el Apéndice~\ref{ap:encuesta}; aquí se describen sus decisiones de diseño.

La sección sociodemográfica (ítems D1--D7) recoge edad, entidad de residencia,
tipo de localidad, nivel educativo propio y situación de actividad. Para
aproximar el nivel socioeconómico de origen se emplearon dos variables
sustitutas (\emph{proxies}): la escolaridad máxima de la madre o tutor
principal (D6) y una autoevaluación retrospectiva de la suficiencia económica
del hogar de crianza (D7). Se optó por estas aproximaciones, de uso común en
encuestas sociales, por ser menos invasivas y menos propensas a la no
respuesta que la pregunta directa por ingreso
[DECIDIR: si el asesor prefiere un instrumento estandarizado, sustituir por la
regla AMAI y ajustar esta justificación].

La sección de percepciones consta de doce ítems (L1--L12) con escala Likert de
cinco puntos, etiquetada en todos sus niveles (de ``totalmente en desacuerdo''
a ``totalmente de acuerdo''), organizados en cuatro bloques: discriminación
lingüística en general (L1--L4), tecnología y trato diferenciado (L5--L6),
escenarios (L7--L9) y experiencia directa (L10--L12). Tres decisiones de
diseño merecen justificación explícita.

Primera, dos ítems (L5 y L7) se redactaron en sentido inverso al resto, de
modo que el acuerdo con ellos expresa la posición contraria a la del bloque.
Su propósito es controlar el sesgo de aquiescencia ---la tendencia a responder
de acuerdo con independencia del contenido--- y detectar patrones de respuesta
automática; ambos se recodifican antes del análisis.

Segunda, el bloque de escenarios sustituye preguntas directas por
mini-viñetas: situaciones concretas (por ejemplo, dos currículums con la misma
experiencia y distinto estilo ante un filtro automatizado) sobre las que el
participante toma postura. Las viñetas reducen la deseabilidad social
---resulta menos costoso emitir un juicio sobre una situación de terceros que
admitir una actitud propia--- y acercan el objeto de la pregunta al escenario
de uso que motiva esta tesis.

Tercera, las cuatro preguntas abiertas (A1--A4) se ordenaron de lo concreto a
lo abstracto: una experiencia vivida de trato diferenciado (A1), las
adaptaciones del propio estilo de escritura ante sistemas automatizados (A2),
el léxico y las formas de escribir que el participante percibe como
penalizadas por una empresa o sistema (A3) y, por último, un mensaje libre a
quienes diseñan estos sistemas (A4). El orden busca anclar primero la memoria
episódica del participante antes de pedirle generalizaciones. El ítem A3
cumple además una función estructural en el diseño: sus respuestas alimentan
directamente la tipología de variaciones del \emph{benchmark}.

Previo al levantamiento, el instrumento se sometió a una prueba piloto con
entre cinco y ocho personas del rango de edad objetivo, con el fin de medir el
tiempo real de respuesta y detectar ítems ambiguos
[PENDIENTE: reportar resultados del piloto y ajustes derivados].

La Tabla~\ref{tab:instrumento-analisis} resume la correspondencia entre las
secciones del instrumento y su uso analítico.

\begin{table}[H]
  \centering
  \caption[Correspondencia entre instrumento y plan de análisis.]{Correspondencia
  entre las secciones del instrumento y su uso en el análisis.}
  \label{tab:instrumento-analisis}
  \begin{tabular}{llp{6.2cm}}
    \toprule
    \textbf{Ítems} & \textbf{Constructo} & \textbf{Uso analítico} \\
    \midrule
    D1--D7   & Perfil sociodemográfico & Caracterización de la muestra \\
    D6--D7   & Nivel socioeconómico (proxy) & Variable de cruce para los bloques Likert \\
    L1--L4   & Discriminación lingüística percibida & Descriptivos y comparación por grupos \\
    L5--L9   & Discriminación tecnológica percibida & Descriptivos y comparación por grupos \\
    L10--L12 & Exposición y confianza & Descriptivos; contexto interpretativo \\
    A1--A2   & Experiencias narradas & Corpus: temas y sentimiento \\
    A3       & Léxico percibido como penalizado & Insumo de la tipología del \emph{benchmark} \\
    A4       & Expectativas normativas & Insumo para la discusión (Cap.~\ref{cap:discusion}) \\
    \bottomrule
  \end{tabular}
\end{table}

\subsection{Análisis cuantitativo}
\label{sec:analisis-cuantitativo}

Las respuestas a los ítems Likert se trataron como variables ordinales. En
consecuencia, el análisis descriptivo reporta distribuciones de frecuencia,
medianas y rangos intercuartílicos por ítem y por bloque, en lugar de medias y
desviaciones estándar, cuya interpretación presupone una escala de intervalo.
Antes del análisis se recodificaron los ítems invertidos (L5 y L7) y se
verificó la consistencia interna de cada bloque mediante el coeficiente alfa
de Cronbach, reportado con carácter descriptivo dada la brevedad de los
bloques [PENDIENTE: reportar valores].

Para explorar si las percepciones varían según el origen socioeconómico, las
respuestas se compararon entre los grupos definidos por las variables proxy
(D6 y D7) mediante pruebas no paramétricas: U de Mann--Whitney para
comparaciones entre dos grupos y H de Kruskal--Wallis para más de dos, con un
nivel de significancia de $\alpha = 0.05$. Dado el carácter no probabilístico
de la muestra, estos contrastes se interpretan como descriptivos de la muestra
obtenida y no como pruebas de hipótesis poblacionales; se reportan junto con
sus tamaños de efecto para facilitar la lectura de la magnitud de las
diferencias [PENDIENTE: completar con los resultados].

\subsection{Análisis del corpus textual}
\label{sec:analisis-corpus}

Las respuestas a las preguntas abiertas se estructuraron en un corpus textual
sobre el que se aplicó la siguiente canalización de procesamiento, implementada
en Python [DECIDIR: fijar versiones de bibliotecas al ejecutar]:

\begin{enumerate}
  \item \textbf{Limpieza y normalización.} Tokenización y lematización con
    \texttt{spaCy} (modelo \texttt{es\_core\_news}), eliminación de palabras
    funcionales y normalización a minúsculas. Una decisión metodológica
    relevante: la normalización se aplicó únicamente sobre la copia de trabajo
    destinada al análisis automático; el texto original de cada respuesta se
    preservó íntegro ---incluidas sus variantes ortográficas---, tanto porque
    esas variantes pueden ser informativas para el objeto de estudio como para
    permitir la cita textual fiel.
  \item \textbf{Análisis de frecuencia.} Frecuencias absolutas y relativas de
    términos y bigramas por pregunta, y ponderación TF--IDF para identificar
    el vocabulario distintivo de cada una.
  \item \textbf{Modelado de temas.} Identificación de temas recurrentes
    mediante asignación latente de Dirichlet (LDA) y, de manera
    complementaria, \texttt{BERTopic} sobre representaciones contextuales en
    español. El número de temas se seleccionó mediante métricas de coherencia
    y revisión manual de interpretabilidad
    [PENDIENTE: reportar configuración y temas resultantes].
  \item \textbf{Análisis de sentimiento.} Clasificación de polaridad con un
    modelo preentrenado para español [DECIDIR: candidato:
    \texttt{pysentimiento}; confirmar al ejecutar], aplicada a las respuestas
    de experiencias (A1--A2) para caracterizar la carga afectiva de las
    narraciones.
\end{enumerate}

El producto final de este análisis es doble: una caracterización temática y
afectiva de las percepciones (reportada en la
Sección~\ref{sec:resultados-encuesta}) y, a partir del ítem A3, un inventario
de rasgos lingüísticos percibidos como penalizados, que constituye el punto de
partida empírico de la tipología de variaciones del \emph{benchmark}
(Sección~\ref{sec:benchmark-diseno}).

\section{Diseño del conjunto de textos controlados}
\label{sec:benchmark-diseno}

\subsection{Principio de control}

El \emph{benchmark} se construye respetando el siguiente principio: dos textos
$t$ y $t'$ del conjunto se consideran un \emph{par controlado} si y sólo si su
contenido informativo es operativamente equivalente y difieren únicamente en
el registro lingüístico. La equivalencia informativa se verifica mediante
revisión humana y, complementariamente, mediante una medida de similitud
semántica sobre representaciones vectoriales.

\subsection{Tipos de variación lingüística considerados}

[Tipología de variaciones a operacionalizar, derivada del análisis de la
encuesta y de la literatura: léxico formal vs.\ coloquial, presencia de
errores ortográficos, uso de anglicismos, longitud oracional, marcadores
discursivos, etc.]

\subsection{Procedimiento de construcción}

[Pasos para construir los pares controlados: redacción inicial, reescritura
con guía estilística, validación por jueces, control de calidad.]

\subsection{Tamaño y composición del \textit{benchmark}}

[Cantidad de textos, número de pares por tipo de variación, distribución de
etiquetas objetivo (en caso de tarea supervisada).]

\section{Evaluación experimental}
\label{sec:experimentos}

\subsection{Modelos evaluados}

[Selección de modelos: al menos un clasificador lineal sobre TF--IDF, al
menos un modelo basado en \textit{transformers} preentrenado en español
(\textit{p.\ ej.,} BETO, RoBERTa en español), y opcionalmente un modelo de
lenguaje grande utilizado como clasificador \textit{zero-shot}.]

\subsection{Tarea de clasificación}

[Definición exacta de la tarea: etiquetas, definiciones operativas, ejemplos.]

\subsection{Protocolo experimental}

\begin{enumerate}
  \item Para cada modelo $f$ y cada par controlado $(t, t')$, se obtienen las
    predicciones $f(t)$ y $f(t')$ y, cuando es posible, las probabilidades
    asociadas $p_f(\cdot \mid t)$ y $p_f(\cdot \mid t')$.
  \item Se calculan métricas estándar de comportamiento agregado por
    \emph{condición de registro}.
  \item Se calcula el \INS según la definición del Capítulo~\ref{cap:ins}.
\end{enumerate}

\subsection{Reproducibilidad}

[Política de versiones de modelos y bibliotecas; semillas aleatorias;
disposición del código y del \emph{benchmark}.]

\section{Consideraciones éticas}
\label{sec:etica}

Las decisiones éticas de esta investigación se concentran en dos frentes: el
tratamiento de las personas participantes en la encuesta y el uso responsable
de los artefactos producidos.

Respecto a las personas participantes, la encuesta se aplicó únicamente a
mayores de edad, con participación voluntaria y posibilidad de abandono en
cualquier momento sin consecuencia alguna. El cuestionario abre con un
consentimiento informado que declara el propósito académico del estudio, la
institución y los responsables, y las condiciones de uso de las respuestas; el
formulario no puede completarse sin aceptarlo expresamente. El instrumento no
solicita nombre, correo, teléfono ni ningún otro dato que permita la
identificación directa de la persona, y las variables sociodemográficas se
recogen en categorías agregadas. Los resultados cuantitativos se reportan
exclusivamente de forma agregada. Las respuestas abiertas pueden citarse
textualmente en este documento, sin atribución; antes de citar cualquier
fragmento se verificó que no contuviera detalles que permitieran la
reidentificación indirecta (nombres de terceros, empresas o lugares
específicos), suprimiéndolos cuando fue necesario
[PENDIENTE: confirmar tras la selección final de citas]. Los datos crudos se
almacenaron en la plataforma del formulario y en una copia local de acceso
restringido al investigador y su asesor, y se conservarán únicamente durante
el desarrollo de la tesis [DECIDIR: confirmar con el asesor si la Facultad
requiere aval de algún comité o registro adicional para encuestas anónimas de
bajo riesgo].

Respecto a los artefactos, el \emph{benchmark} de pares controlados merece una
consideración específica. Por construcción, contiene textos redactados en
registros socialmente estigmatizados; su propósito es auditar sistemas de
clasificación, pero un conjunto así podría, en principio, emplearse en sentido
contrario ---por ejemplo, para entrenar deliberadamente un clasificador que
penalice esos registros. Para reducir ese riesgo, el \emph{benchmark} se
documenta y distribuye con una declaración explícita de uso previsto
(auditoría y evaluación de equidad), siguiendo la práctica de documentación de
conjuntos de datos discutida en la literatura, y los textos que lo integran
son simulados: ningún par proviene de currículums ni de textos reales de
personas. Finalmente, los experimentos del Capítulo~\ref{cap:resultados}
evalúan modelos preentrenados disponibles públicamente y no implican
recolección adicional de datos personales.

% =============================================================================
%  Capítulo 4 -- Resultados
% =============================================================================
\chapter{Resultados}
\label{cap:resultados}

\section{Análisis de percepciones sobre discriminación tecnológica}
\label{sec:resultados-encuesta}

\subsection{Caracterización de la muestra}

[Tablas y gráficos descriptivos sobre las características demográficas y
socioeconómicas de la muestra obtenida.]

% Plantilla de tabla con booktabs:
\begin{table}[H]
  \centering
  \caption[Caracterización de la muestra.]{Caracterización de la muestra
  (placeholder).}
  \label{tab:muestra}
  \begin{tabular}{lrr}
    \toprule
    \textbf{Variable} & \textbf{Frecuencia} & \textbf{Porcentaje} \\
    \midrule
    {[Categoría 1]} & --- & --- \\
    {[Categoría 2]} & --- & --- \\
    {[Categoría 3]} & --- & --- \\
    \midrule
    Total         & --- & 100.0 \\
    \bottomrule
  \end{tabular}
\end{table}

\subsection{Percepciones cuantitativas}

[Reporte de las respuestas a preguntas Likert.]

\subsection{Análisis del corpus de respuestas abiertas}

[Frecuencias, temas identificados, sentimientos, ejemplos representativos
seleccionados con criterios explícitos.]

\subsection{Síntesis: dimensiones relevantes para el \textit{benchmark}}

[Lista de tipos de variación lingüística que la evidencia recogida sugiere
como prioritarios para el diseño del \emph{benchmark}.]

\section{Evaluación del comportamiento del sistema frente a variaciones lingüísticas}
\label{sec:resultados-experimentos}

\subsection{Descripción del \textit{benchmark} construido}

[Resumen cuantitativo del \emph{benchmark}: número de pares, distribución por
tipo de variación, longitud promedio de los textos, validaciones realizadas.]

\subsection{Resultados agregados por modelo}

[Tabla comparativa de métricas estándar para cada modelo evaluado, separadas
por condición de registro.]

\subsection{Resultados por tipo de variación lingüística}

[Análisis desagregado: ¿qué tipos de variación generan mayor inconsistencia?]

\subsection{Casos cualitativos seleccionados}

[Ejemplos ilustrativos de pares en los que el modelo cambia su predicción
ante variación de registro, presentados con discusión.]

\section{Identificación de patrones de sesgo sociolectal}
\label{sec:patrones}

[Síntesis de los hallazgos cuantitativos y cualitativos; identificación de
patrones sistemáticos; preparación para el cálculo del \INS en el siguiente
capítulo.]

% =============================================================================
%  Capítulo 5 -- Propuesta del Índice de Neutralidad Sociolectal
% =============================================================================
\chapter{Propuesta del Índice de Neutralidad Sociolectal}
\label{cap:ins}

\section{Definición conceptual del índice}
\label{sec:ins-conceptual}

El \emph{Índice de Neutralidad Sociolectal} (\INS) busca operacionalizar una
intuición sencilla: un sistema de clasificación de texto es tanto más
\emph{neutral} respecto a la variación sociolectal cuanto más consistente sea
su comportamiento al ser confrontado con pares de textos que comparten
contenido informativo pero difieren en registro lingüístico.

El \INS no es una métrica de \emph{exactitud}, sino de \emph{consistencia
intra-par}. Un modelo puede tener exactitud baja y, sin embargo, \INS alto si
sus errores se distribuyen de manera independiente del registro; y un modelo
puede tener exactitud alta y, no obstante, \INS bajo si su comportamiento
correcto se concentra sistemáticamente en un registro y no en otro.

\section{Formalización preliminar}
\label{sec:ins-formal}

\begin{definicion}[Par controlado]
\label{def:par-controlado}
Un \emph{par controlado} es un par ordenado $(t, t')$ de textos tales que su
contenido informativo, evaluado bajo los criterios de la
Sección~\ref{sec:benchmark-diseno}, es operativamente equivalente, y que
difieren únicamente en su registro lingüístico.
\end{definicion}

Sea $\mathcal{B} = \{(t_i, t_i')\}_{i=1}^{N}$ un \emph{benchmark} de pares
controlados y sea $f$ un clasificador. Se proponen dos formulaciones
complementarias del \INS, según la información de salida disponible.

\subsection{Formulación basada en decisiones}

Cuando solo se dispone de la decisión $f(t) \in \mathcal{Y}$:

\begin{definicion}[\INS de decisión]
\label{def:ins-decision}
\[
  \mathrm{INS}_{\mathrm{dec}}(f, \mathcal{B}) \;=\;
  \frac{1}{N} \sum_{i=1}^{N} \mathbb{1}\!\left[ f(t_i) = f(t_i') \right].
\]
\end{definicion}

\noindent
Esta formulación corresponde a la proporción de pares en los que el modelo
emite la misma decisión para ambos miembros. Toma valores en $[0,1]$, con $1$
indicando consistencia perfecta de decisión.

\subsection{Formulación basada en probabilidades}

Cuando se dispone de la distribución posterior $p_f(\cdot \mid t)$ sobre las
clases:

\begin{definicion}[\INS probabilístico]
\label{def:ins-prob}
\[
  \mathrm{INS}_{\mathrm{prob}}(f, \mathcal{B}) \;=\;
  1 \;-\; \frac{1}{N} \sum_{i=1}^{N}
  d\!\left( p_f(\cdot \mid t_i),\; p_f(\cdot \mid t_i') \right),
\]
\noindent donde $d(\cdot,\cdot)$ es una medida de divergencia entre
distribuciones (por ejemplo, divergencia de Jensen--Shannon o distancia de
variación total), normalizada al intervalo $[0,1]$.
\end{definicion}

\noindent
La formulación probabilística captura cambios en la confianza del modelo que
la formulación de decisión ignora, y por tanto es preferible cuando se dispone
de los puntajes de probabilidad.

\subsection{Propiedades deseables}

\begin{proposicion}[Acotación]
$\mathrm{INS}_{\mathrm{dec}}(f, \mathcal{B}), \mathrm{INS}_{\mathrm{prob}}(f, \mathcal{B}) \in [0,1]$.
\end{proposicion}

\begin{proposicion}[Invariancia a permutaciones del \textit{benchmark}]
El valor de ambas formulaciones es invariante a cualquier reordenamiento de
los pares de $\mathcal{B}$.
\end{proposicion}

[Discutir y, según corresponda, demostrar otras propiedades: comportamiento
ante \emph{benchmarks} desbalanceados por tipo de variación, robustez a
ruido, etc.]

\subsection{Variantes desagregadas}

Si los pares de $\mathcal{B}$ están etiquetados por tipo de variación
lingüística $v \in V$, se define el \INS condicional como

\[
  \mathrm{INS}_{\mathrm{dec}}(f, \mathcal{B} \mid v)
  \;=\;
  \frac{1}{|\mathcal{B}_v|}
  \sum_{(t,t') \in \mathcal{B}_v}
  \mathbb{1}\!\left[ f(t) = f(t') \right],
\]
\noindent donde $\mathcal{B}_v$ es el subconjunto de pares cuyo tipo de
variación es $v$. Esta variante permite identificar qué tipos de variación
afectan más al modelo.

\section{Procedimiento de cálculo}
\label{sec:ins-procedimiento}

Esta sección operacionaliza las definiciones de la
Sección~\ref{sec:ins-formal} en un procedimiento ejecutable y reproducible.
Se describen las entradas requeridas, el algoritmo paso a paso, las
decisiones de implementación y una implementación de referencia.

\subsection{Entradas y supuestos}

El cálculo del \INS requiere dos entradas. La primera es un \emph{benchmark}
$\mathcal{B} = \{(t_i, t_i')\}_{i=1}^{N}$ de pares controlados
(Definición~\ref{def:par-controlado}), donde cada par está opcionalmente
etiquetado con su tipo de variación lingüística $v \in V$ para permitir la
desagregación. La segunda es un clasificador $f$ que expone, como mínimo, una
interfaz de decisión $f(t) \in \mathcal{Y}$ y, de manera opcional, la
distribución de probabilidad posterior $p_f(\cdot \mid t)$ sobre las clases.

El procedimiento presupone que las salidas del modelo son deterministas o han
sido vueltas reproducibles. Para modelos con componentes estocásticos ---en
particular, modelos de lenguaje grandes usados como clasificadores--- esto
exige fijar la temperatura de muestreo en cero o, cuando ello no sea posible,
promediar las salidas sobre $k$ ejecuciones independientes y documentar $k$.
Sin esta precaución, la inconsistencia atribuible al azar del muestreo se
confundiría con la inconsistencia sociolectal que el índice pretende medir.

\subsection{Algoritmo}

El cálculo procede en cinco pasos.

\begin{enumerate}
  \item \textbf{Verificación del \emph{benchmark}.} Se comprueba que
    $\mathcal{B}$ no esté vacío, que cada par esté completo y que las
    etiquetas de tipo de variación, si existen, pertenezcan al catálogo $V$
    definido en la Sección~\ref{sec:benchmark-diseno}. Un \emph{benchmark}
    vacío no produce un \INS de cero sino un valor \emph{indefinido}: el
    procedimiento debe fallar explícitamente en ese caso, no devolver un
    número que pueda malinterpretarse.
  \item \textbf{Obtención de salidas.} Para cada par $(t_i, t_i')$ se
    obtienen $f(t_i)$ y $f(t_i')$ y, si el modelo lo permite,
    $p_f(\cdot \mid t_i)$ y $p_f(\cdot \mid t_i')$. Esto implica exactamente
    $2N$ invocaciones al modelo; el costo del procedimiento es lineal en el
    tamaño del \emph{benchmark}.
  \item \textbf{Cálculo de la formulación de decisión.} Se computa
    $\mathrm{INS}_{\mathrm{dec}}$ según la Definición~\ref{def:ins-decision}:
    la proporción de pares con decisión coincidente.
  \item \textbf{Cálculo de la formulación probabilística.} Si se dispone de
    probabilidades, se computa $\mathrm{INS}_{\mathrm{prob}}$ según la
    Definición~\ref{def:ins-prob}, con la divergencia de Jensen--Shannon como
    medida $d$ (la elección se justifica más abajo).
  \item \textbf{Desagregación y reporte.} Se calcula el \INS condicional por
    tipo de variación, $\mathrm{INS}(f, \mathcal{B} \mid v)$, y se reporta el
    índice global junto con el perfil desagregado, el tamaño $N$ del
    \emph{benchmark}, el tamaño $N_v$ de cada celda y las métricas de
    desempeño del modelo sobre la tarea. La razón de este último requisito se
    desarrolla en la Sección~\ref{sec:ins-interpretacion}.
\end{enumerate}

\subsection{Decisiones de implementación}

\paragraph{Elección y normalización de la divergencia.} La
Definición~\ref{def:ins-prob} exige una divergencia $d$ acotada en $[0,1]$.
Se adopta la divergencia de Jensen--Shannon calculada con logaritmo en base 2,
que satisface esta cota de manera natural ---alcanza $0$ para distribuciones
idénticas y $1$ para distribuciones con soportes disjuntos--- sin requerir
normalización \emph{ad hoc}. Es además simétrica y está definida incluso
cuando alguna clase recibe probabilidad cero, pues la mezcla
$m = \tfrac{1}{2}(p + p')$ es estrictamente positiva dondequiera que alguna
de las dos distribuciones lo sea.

\paragraph{Degradación controlada.} Cuando el modelo no expone
probabilidades ---caso frecuente en servicios de clasificación accedidos como
caja negra--- el procedimiento se degrada a la formulación de decisión y el
reporte lo hace explícito. Comparar un $\mathrm{INS}_{\mathrm{prob}}$ de un
modelo con un $\mathrm{INS}_{\mathrm{dec}}$ de otro no es válido: son
cantidades distintas.

\paragraph{Separación entre evaluación y métrica.} La implementación separa
la obtención de salidas del modelo (paso 2) del cálculo del índice (pasos
3--5). Esta separación permite probar unitariamente las funciones de cálculo
con casos sintéticos de respuesta conocida y recalcular el índice sin repetir
las $2N$ invocaciones al modelo, que son la parte costosa del procedimiento.

\subsection{Implementación de referencia}

El Código~\ref{lst:ins-ref} presenta la implementación de referencia en
Python, sin más dependencia que \texttt{NumPy}. Las funciones de cálculo
fueron verificadas con pruebas unitarias que cubren los casos extremos:
consistencia perfecta ($\INS = 1$), inversión total ($\INS = 0$), cotas de la
divergencia y rechazo del \emph{benchmark} vacío.

\begin{lstlisting}[caption={Implementación de referencia del cálculo del \INS.},label={lst:ins-ref}]
import numpy as np

def divergencia_js(p, q, base=2.0):
    """Divergencia de Jensen-Shannon entre dos distribuciones
    discretas. Con base 2, el resultado queda acotado en [0, 1]."""
    p = np.asarray(p, dtype=float); q = np.asarray(q, dtype=float)
    p = p / p.sum(); q = q / q.sum()
    m = 0.5 * (p + q)
    def kl(a, b):
        mascara = a > 0   # m > 0 donde a > 0: el cociente es seguro
        return float(np.sum(a[mascara]
                     * (np.log(a[mascara] / b[mascara]) / np.log(base))))
    return 0.5 * kl(p, m) + 0.5 * kl(q, m)

def calcular_ins_dec(pares):
    """INS de decision. `pares` es una lista de tuplas (y, y_prima)
    con las decisiones del modelo para cada par controlado."""
    if not pares:
        raise ValueError("Benchmark vacio: el INS no esta definido.")
    return sum(int(y == yp) for y, yp in pares) / len(pares)

def calcular_ins_prob(pares, base=2.0):
    """INS probabilistico. `pares` es una lista de tuplas (p, p_prima)
    con las distribuciones posteriores del modelo para cada par."""
    if not pares:
        raise ValueError("Benchmark vacio: el INS no esta definido.")
    total = sum(divergencia_js(p, q, base) for p, q in pares)
    return 1.0 - total / len(pares)

def ins_por_variacion(registros):
    """INS de decision desagregado por tipo de variacion.
    `registros` es un iterable de tuplas (tipo, y, y_prima)."""
    grupos = {}
    for tipo, y, yp in registros:
        grupos.setdefault(tipo, []).append((y, yp))
    return {t: calcular_ins_dec(ps) for t, ps in grupos.items()}
\end{lstlisting}

\subsection{Ejemplo ilustrativo}

Para fijar ideas, considérese un \emph{benchmark} ficticio de cuatro pares
sobre una tarea binaria con etiquetas \{apto, no apto\}, construido con fines
puramente ilustrativos ---los valores no provienen de ningún experimento de
esta tesis. La Tabla~\ref{tab:ejemplo-ins} muestra las decisiones de un
modelo hipotético sobre cada miembro del par.

\begin{table}[H]
  \centering
  \caption[Ejemplo ilustrativo del cálculo del \INS de decisión.]{Ejemplo
  ilustrativo (datos ficticios) del cálculo de
  $\mathrm{INS}_{\mathrm{dec}}$.}
  \label{tab:ejemplo-ins}
  \begin{tabular}{clll}
    \toprule
    \textbf{Par} & $f(t_i)$ & $f(t_i')$ & \textbf{¿Coincide?} \\
    \midrule
    1 & apto    & apto    & sí \\
    2 & apto    & no apto & no \\
    3 & no apto & no apto & sí \\
    4 & apto    & apto    & sí \\
    \bottomrule
  \end{tabular}
\end{table}

Tres de los cuatro pares reciben la misma decisión, de modo que
$\mathrm{INS}_{\mathrm{dec}} = 3/4 = 0.75$. El ejemplo permite además
ilustrar qué añade la formulación probabilística: supóngase que en el par~1
el modelo asigna a la clase ``apto'' una probabilidad de $0.91$ para $t_1$
pero de $0.62$ para $t_1'$. La decisión no cambia ---ambas superan el umbral
de $0.5$--- y el par cuenta como consistente para
$\mathrm{INS}_{\mathrm{dec}}$; sin embargo, la divergencia de
Jensen--Shannon entre $(0.91, 0.09)$ y $(0.62, 0.38)$ es de
aproximadamente $0.089$, un desplazamiento de confianza que
$\mathrm{INS}_{\mathrm{prob}}$ sí registra. En tareas donde la probabilidad
alimenta decisiones posteriores ---ordenar candidatos en lugar de solo
aceptarlos o rechazarlos--- ese desplazamiento tiene consecuencias aunque la
etiqueta no cambie.

\section{Aplicación a los resultados experimentales}
\label{sec:ins-aplicacion}

[Reporte de los valores de \INS calculados sobre los modelos del
Capítulo~\ref{cap:resultados}, en sus formulaciones agregada y desagregada
por tipo de variación, con su correspondiente interpretación.]

\section{Interpretación y alcances del indicador}
\label{sec:ins-interpretacion}

Un indicador solo es útil si se conocen con precisión sus condiciones de
lectura. Esta sección delimita qué afirma y qué no afirma un valor del \INS,
bajo qué condiciones es comparable y cómo se sitúa frente a las métricas de
equidad algorítmica existentes.

\subsection{Qué mide y qué no mide el \INS}

El \INS mide \emph{consistencia intra-par}: el grado en que un clasificador
responde igual ante dos textos cuyo contenido es equivalente y cuyo registro
difiere. Cuatro precisiones acotan esta lectura.

Primera, consistencia no es corrección. Ambas dimensiones son ortogonales y
su cruce genera cuatro escenarios cualitativamente distintos: un modelo
exacto y consistente (el caso deseable); un modelo exacto en el agregado pero
inconsistente entre registros, que es precisamente el sesgo que las métricas
agregadas encubren; un modelo inexacto pero consistente, que falla por igual
para todos los registros; y un modelo inexacto e inconsistente. El \INS
distingue las columnas de este cruce, no las filas: por eso debe reportarse
siempre junto a las métricas de desempeño de la tarea, nunca en su lugar.

Segunda, y como caso límite de lo anterior: la neutralidad perfecta es
trivialmente alcanzable. Un clasificador constante, que ignora por completo
su entrada, obtiene $\mathrm{INS}_{\mathrm{dec}} = 1$ por construcción. Un
\INS alto, aislado de toda métrica de desempeño, no acredita nada: el índice
es informativo únicamente para modelos que resuelven la tarea con un
desempeño no trivial.

Tercera, el \INS es simétrico y, por tanto, no direccional: registra
\emph{que} el modelo cambia su respuesta ante el registro, pero no
\emph{a quién} favorece el cambio. Un mismo valor de \INS es compatible con
un modelo que penaliza sistemáticamente el registro coloquial y con uno que
lo favorece. Establecer la dirección del sesgo exige un análisis
complementario ---por ejemplo, comparar la tasa de asignación de la clase
favorable por condición de registro---, que en esta tesis se reporta en la
Sección~\ref{sec:resultados-experimentos}.

Cuarta, el \INS es una medida relativa al \emph{benchmark}, no una propiedad
intrínseca del modelo. Su validez hereda la validez de la operacionalización
de ``equivalencia de contenido'' descrita en la
Sección~\ref{sec:benchmark-diseno}: si los pares confunden variación de
registro con variación de contenido, el índice atribuirá al sesgo
sociolectal inconsistencias que en realidad responden a información
distinta. Todo valor de \INS debe leerse, en rigor, como
``\INS del modelo $f$ \emph{sobre el benchmark} $\mathcal{B}$''.

Finalmente, conviene precisar la relación entre las dos formulaciones. La
probabilística es más sensible: detecta desplazamientos de confianza que no
alcanzan a cruzar el umbral de decisión, como ilustró el ejemplo de la
Sección~\ref{sec:ins-procedimiento}. La de decisión, en cambio, refleja la
consecuencia material: lo que la persona evaluada experimenta es la decisión,
no la distribución interna del modelo. Reportar ambas, cuando es posible,
ofrece la lectura más completa: la probabilística como instrumento de
diagnóstico fino, la de decisión como medida del efecto práctico.

\subsection{Umbrales y comparabilidad}

No existe un umbral universal a partir del cual un valor de \INS certifique
neutralidad, y esta tesis deliberadamente no propone uno. Lo que sí pueden
ofrecerse son puntos de referencia que delimitan el rango interpretable. En
una tarea binaria, un clasificador que respondiera al azar de manera
independiente sobre cada texto obtendría un
$\mathrm{INS}_{\mathrm{dec}}$ esperado de $0.5$; el clasificador constante,
como se señaló, alcanza $1$ de forma trivial. Los valores informativos de un
modelo real se leen contra ese telón de fondo y, sobre todo, en
\emph{comparación}: entre modelos evaluados sobre el mismo \emph{benchmark},
entre versiones sucesivas de un mismo modelo, o entre tipos de variación
dentro de un mismo modelo.

De ahí se siguen las condiciones de comparabilidad. Dos valores de \INS son
directamente comparables solo si provienen del mismo \emph{benchmark} ---los
mismos pares, la misma tarea, el mismo espacio de etiquetas--- y de la misma
formulación (decisión con decisión, probabilística con probabilística).
Comparar valores obtenidos sobre \emph{benchmarks} distintos no es válido sin
más: un conjunto de pares con variaciones sutiles producirá índices
sistemáticamente más altos que uno con variaciones marcadas, sin que ello
diga nada sobre los modelos.

Por último, el perfil desagregado
$\{\mathrm{INS}(f, \mathcal{B} \mid v)\}_{v \in V}$ es más informativo que el
escalar global, pues localiza el problema: un modelo puede ser robusto ante
anglicismos e inconsistente ante variación ortográfica, y esa estructura
---invisible en el agregado--- es la que orienta tanto el diagnóstico como
una eventual mitigación. El protocolo de reporte adoptado en esta tesis
comprende, por ello, el \INS global en sus dos formulaciones cuando es
posible, el perfil desagregado por tipo de variación con el tamaño $N_v$ de
cada celda, y las métricas de desempeño del modelo sobre la tarea.

\subsection{Relación con otras métricas de equidad}

La literatura sobre equidad algorítmica ha producido una familia amplia de
criterios formales, sistematizada en revisiones como la de Mehrabi y
colaboradores \cite{Mehrabi2021Survey}. Situar el \INS en ese mapa ayuda a
precisar su aporte.

Las métricas de \emph{equidad de grupo} ---la paridad demográfica, que exige
tasas de asignación de la clase favorable iguales entre grupos, o la igualdad
de oportunidades, que exige tasas de verdaderos positivos iguales
condicionadas a la etiqueta real--- comparan agregados calculados sobre
individuos \emph{distintos} pertenecientes a grupos distintos
\cite{Mehrabi2021Survey}. El \INS opera en otra unidad de análisis: compara
las respuestas del modelo ante versiones del \emph{mismo} contenido, lo que
lo aproxima a la familia de la \emph{equidad individual}, cuyo principio es
que entradas similares reciban tratamientos similares. El \emph{benchmark} de
pares controlados es, en efecto, una operacionalización textual de
``similitud'': dos textos se declaran similares cuando su contenido es
equivalente y solo difiere su registro.

Esta construcción acerca el \INS, en particular, a la noción de
\emph{equidad contrafáctica}, según la cual una decisión es justa si no
cambiaría en el mundo contrafáctico donde el individuo perteneciera a otro
grupo protegido [FUENTE: añadir la referencia primaria de equidad
contrafáctica (Kusner et al., 2017, \emph{Counterfactual Fairness}, NeurIPS)
a referencias.bib; mientras tanto, la noción está cubierta por el survey de
Mehrabi]. Cada par $(t, t')$ del \emph{benchmark} funciona como un
contrafáctico textual: el ``individuo'' es el contenido informativo y la
intervención consiste en cambiar únicamente el atributo sensible, el
registro. La diferencia importante es de estatus epistémico: la equidad
contrafáctica plantea la intervención sobre un modelo causal del mundo,
mientras que el \INS la realiza materialmente sobre el texto, que es donde el
atributo sociolectal se manifiesta. Esa materialización es lo que vuelve al
índice computable sin supuestos causales adicionales.

El \INS comparte también una afinidad metodológica con las pruebas de
invariancia usadas en la evaluación conductual de modelos de \PLN, en las que
se verifica que perturbaciones que no alteran el significado no alteren la
predicción [FUENTE: añadir la referencia de evaluación conductual por
invariancia (Ribeiro et al., 2020, \emph{CheckList}, ACL) a
referencias.bib]. Bajo esa lectura, el \INS es una prueba de invariancia cuya
perturbación no es sintética ni arbitraria, sino sociolingüísticamente
anclada: las variaciones provienen de los marcadores de registro
documentados en el Capítulo~\ref{cap:marco} y derivados empíricamente de la
encuesta.

Dos ventajas prácticas se siguen de esta posición en el mapa. Primera, a
diferencia de la igualdad de oportunidades, el \INS en su formulación de
consistencia no requiere etiquetas de verdad de terreno: basta el
\emph{benchmark} y el acceso al modelo, lo que lo hace aplicable a sistemas
de caja negra cuya tarea ni siquiera tiene una etiqueta ``correcta''
verificable. Segunda, su unidad de análisis pareada lo protege de un modo de
fallo conocido de las métricas de grupo: un sistema puede satisfacer la
paridad demográfica compensando errores entre subpoblaciones sin tratar de
manera consistente a ningún individuo en particular. Nada de esto convierte
al \INS en sustituto de aquellas métricas: mide una dimensión ---la
consistencia ante variación sociolectal--- que las métricas de grupo no
capturan, y omite otras ---las tasas de acierto condicionadas a la verdad de
terreno--- que ellas sí. Su lugar es el de un complemento dentro de una
batería de evaluación, no el de una medida única de equidad.

% =============================================================================
%  Capítulo 6 -- Discusión
% =============================================================================
\chapter{Discusión}
\label{cap:discusion}

\section{Implicaciones sociotécnicas de los resultados}
\label{sec:implicaciones}

[Articular los hallazgos cuantitativos del Capítulo~\ref{cap:resultados} y
los valores del \INS del Capítulo~\ref{cap:ins} con la discusión teórica del
Capítulo~\ref{cap:marco}. ¿Qué dicen estos resultados sobre la traslación de
la discriminación lingüística humana al ámbito algorítmico? ¿Qué dicen sobre
las prácticas actuales de evaluación de sistemas de \PLN desplegados en
contextos hispanohablantes?]

\subsection{Implicaciones para el desarrollo de sistemas}

[Lecciones aplicables a equipos que diseñan y despliegan clasificadores de
texto en contextos como el filtrado de currículums.]

\subsection{Implicaciones para la evaluación de equidad}

[¿Cuándo conviene incorporar el \INS o variantes suyas en las baterías de
evaluación de un sistema de \PLN?]

\subsection{Implicaciones para la discusión pública sobre algoritmos}

[Relación de los resultados con debates más amplios sobre transparencia,
rendición de cuentas y regulación de sistemas algorítmicos.]

\section{Limitaciones metodológicas del estudio}
\label{sec:limitaciones}

\subsection{Limitaciones del componente social}

[Sobre la muestra de la encuesta, su representatividad, los sesgos del
instrumento.]

\subsection{Limitaciones del \textit{benchmark} controlado}

[Naturaleza simulada de los textos; cobertura limitada de los tipos de
variación; subjetividad en el control de equivalencia informativa.]

\subsection{Limitaciones del indicador propuesto}

[El \INS como indicador preliminar; supuestos implícitos; condiciones bajo
las cuales podría sobreestimar o subestimar la neutralidad efectiva del
sistema.]

\subsection{Limitaciones de la generalización}

[Acotaciones explícitas sobre la transferibilidad de los hallazgos a otras
lenguas, contextos y dominios de aplicación.]

% =============================================================================
%  Capítulo 7 -- Conclusiones y trabajo futuro
% =============================================================================
\chapter{Conclusiones y trabajo futuro}
\label{cap:conclusiones}

\section{Conclusiones generales}
\label{sec:conclusiones-generales}

[Síntesis de los hallazgos principales de la tesis, organizada en torno a los
cuatro objetivos específicos planteados en la Sección~\ref{sec:objetivos}.]

\begin{enumerate}[label=\textbf{C\arabic*.},leftmargin=2.2em]
  \item \textbf{Respecto al objetivo O1} (documentación de clasismo
    estructural y discriminación lingüística en el contexto mexicano):
    [síntesis].
  \item \textbf{Respecto al objetivo O2} (diseño del \emph{benchmark} de
    textos controlados): [síntesis].
  \item \textbf{Respecto al objetivo O3} (evaluación experimental del
    comportamiento del sistema): [síntesis].
  \item \textbf{Respecto al objetivo O4} (diseño y cálculo del \INS):
    [síntesis].
\end{enumerate}

\section{Líneas de investigación futura}
\label{sec:trabajo-futuro}

\subsection{Validación empírica del \INS}

[Estudios sistemáticos para caracterizar el comportamiento del indicador en
distintas familias de modelos y \emph{benchmarks}, comparación con otras
métricas de equidad, calibración de umbrales orientativos.]

\subsection{Expansión del \textit{benchmark}}

[Incorporación de variantes regionales del español mexicano; cobertura de
tareas distintas a la clasificación binaria; ampliación a otros idiomas.]

\subsection{Mitigación de sesgos sociolectales}

[Aproximaciones de mitigación inspiradas en los hallazgos:
\textit{data augmentation} con paráfrasis controladas, entrenamiento con
restricciones de invariancia, ajustes de \textit{post-processing}.]

\subsection{Articulación con marcos regulatorios}

[Posibilidades de uso del \INS o de indicadores análogos en marcos
regulatorios emergentes sobre inteligencia artificial.]


% =============================================================================
%  Materia posterior (apéndices, bibliografía)
% =============================================================================
\appendix
% =============================================================================
%  Apéndice A -- Instrumento de la encuesta
%  -------------------------------------------------------------------------
%  NOTA PARA AXEL:
%  - Este es un BORRADOR COMPLETO del instrumento, listo para discutir con tu
%    asesor antes de aplicarlo. Tres decisiones de diseño que debes validar:
%    (1) escala Likert de 5 puntos con etiqueta en cada punto;
%    (2) los ítems L7-L9 usan viñetas (mini-escenarios) en vez de preguntas
%        directas, para reducir deseabilidad social;
%    (3) las abiertas A1-A4 están ordenadas de lo concreto a lo abstracto.
%  - Marcas [DECIDIR]: puntos donde la decisión es tuya o de tu asesor.
%  - Antes de aplicar: pilotear con 5-8 personas y medir tiempo de respuesta
%    (objetivo: menos de 12 minutos).
% =============================================================================
\chapter{Instrumento de la encuesta}
\label{ap:encuesta}

\section*{Ficha técnica del instrumento}

\begin{description}
  \item[Título.] Percepciones sobre lenguaje, tecnología y trato diferenciado.
  \item[Población objetivo.] Personas de 18 a 30 años residentes en México.
    [DECIDIR: confirmar rango de edad con el asesor.]
  \item[Modo de aplicación.] Formulario en línea, autoadministrado, anónimo.
  \item[Duración estimada.] 10--12 minutos.
  \item[Estructura.] Cuatro secciones: consentimiento informado, datos
    sociodemográficos, percepciones (escala Likert) y preguntas abiertas.
\end{description}

\section*{Sección 1. Consentimiento informado}

\begin{quote}
Esta encuesta forma parte de una investigación de tesis de licenciatura en
Ciencias de la Computación (UNAM) sobre la relación entre el estilo de
escritura y el trato que las personas reciben de sistemas automatizados
(por ejemplo, filtros de currículums o moderación de contenido en plataformas).

Tu participación es \textbf{voluntaria y anónima}: no se solicita tu nombre ni
ningún dato que permita identificarte. Las respuestas se utilizarán únicamente
con fines académicos y se reportarán de forma agregada; las respuestas a
preguntas abiertas podrán citarse textualmente en la tesis sin atribución.
Puedes abandonar el cuestionario en cualquier momento sin consecuencia alguna.

Responsable: Axel Casas Espinosa (correo de contacto: [DECIDIR]).
Asesor: M.\ en C.I.C.\ Manuel Soto Romero, Facultad de Ciencias, UNAM.

$\square$ He leído lo anterior y acepto participar.
\end{quote}

\section*{Sección 2. Datos sociodemográficos}

\begin{enumerate}[label=\textbf{D\arabic*.}]
  \item Edad (años cumplidos): \rule{2.5cm}{0.4pt}
  \item Estado de residencia: \rule{5cm}{0.4pt}
  \item ¿Tu localidad es principalmente...? \quad
    $\square$ Urbana \quad $\square$ Suburbana \quad $\square$ Rural
  \item Último nivel educativo concluido: \quad
    $\square$ Secundaria o menos \quad $\square$ Bachillerato \quad
    $\square$ Licenciatura \quad $\square$ Posgrado
  \item ¿Actualmente estudias, trabajas, ambas o ninguna? \quad
    $\square$ Estudio \quad $\square$ Trabajo \quad $\square$ Ambas \quad
    $\square$ Ninguna
  \item Máximo nivel educativo concluido por tu madre o tutor(a) principal:
    \quad $\square$ Primaria o menos \quad $\square$ Secundaria \quad
    $\square$ Bachillerato \quad $\square$ Licenciatura o más \quad
    $\square$ No sé
  \item Pensando en el hogar donde creciste, ¿con cuál de estas frases lo
    describirías mejor? \quad
    $\square$ El dinero no alcanzaba ni para lo básico \quad
    $\square$ Alcanzaba justo para lo básico \quad
    $\square$ Alcanzaba para lo básico y algo más \quad
    $\square$ Vivíamos con holgura
\end{enumerate}

\noindent\textit{Nota metodológica: D6 y D7 son aproximaciones (proxies) de
nivel socioeconómico de origen, menos invasivas que preguntar ingreso. Si el
asesor prefiere un instrumento estandarizado, considerar la regla AMAI
[DECIDIR].}

\section*{Sección 3. Percepciones (escala Likert)}

\noindent Indica qué tan de acuerdo estás con cada afirmación: \\
\textbf{1} = Totalmente en desacuerdo \quad \textbf{2} = En desacuerdo \quad
\textbf{3} = Ni de acuerdo ni en desacuerdo \quad \textbf{4} = De acuerdo
\quad \textbf{5} = Totalmente de acuerdo

\subsection*{Bloque I. Discriminación lingüística en general}

\begin{enumerate}[label=\textbf{L\arabic*.}]
  \item En México, la forma de hablar o escribir de una persona influye en
    cómo la tratan los demás.
  \item He visto que se juzgue a alguien como menos inteligente o menos capaz
    por su manera de expresarse.
  \item Yo he modificado mi forma de hablar o escribir para causar una mejor
    impresión en contextos formales (escuela, trabajo, trámites).
  \item Alguna vez me he sentido juzgado(a) o tratado(a) diferente por mi
    forma de hablar o escribir.
\end{enumerate}

\subsection*{Bloque II. Tecnología y trato diferenciado}

\begin{enumerate}[label=\textbf{L\arabic*.},resume]
  \item Creo que los sistemas automatizados (filtros de currículums,
    moderación en redes, atención por chat) tratan igual a todas las personas
    sin importar cómo escriban.\footnote{Ítem invertido: se redacta en sentido
    contrario a los demás para detectar respuesta automática. Recodificar al
    analizar.}
  \item Me preocupa que una computadora evalúe mis textos (por ejemplo, mi
    currículum) sin intervención humana.
\end{enumerate}

\subsection*{Bloque III. Escenarios}

\noindent\textit{Lee cada situación y responde con la misma escala.}

\begin{enumerate}[label=\textbf{L\arabic*.},resume]
  \item Dos personas con la misma experiencia laboral envían su currículum
    a una empresa que usa un filtro automatizado. Una escribe de manera muy
    formal; la otra, de manera sencilla y directa, con algunos modismos.
    \textit{Creo que el filtro las evaluaría igual.}\footnotemark[\value{footnote}]
  \item Si un sistema automatizado rechazara mi solicitud (de empleo, beca o
    crédito), sospecharía que mi forma de escribir influyó en el resultado.
  \item Considero justo que una empresa use la ``calidad de redacción'' como
    criterio automático para filtrar candidatos, aunque el puesto no requiera
    escribir.
\end{enumerate}

\subsection*{Bloque IV. Experiencia directa}

\begin{enumerate}[label=\textbf{L\arabic*.},resume]
  \item He interactuado con sistemas automatizados que evalúan o filtran lo
    que escribo (postulaciones en línea, chatbots de soporte, moderación en
    redes sociales).
  \item En esas interacciones, he sentido que el resultado habría sido
    distinto si hubiera escrito de otra manera.
  \item Confío en que las decisiones automatizadas sobre textos pueden
    apelarse o revisarse por una persona si lo solicito.
\end{enumerate}

\section*{Sección 4. Preguntas abiertas}

\noindent\textit{Responde con tus propias palabras; no hay respuestas
correctas ni incorrectas. La extensión es libre.}

\begin{enumerate}[label=\textbf{A\arabic*.}]
  \item Cuéntanos alguna ocasión en la que sientas que a ti o a alguien
    cercano lo trataron diferente (mejor o peor) por su forma de hablar o
    escribir. ¿Qué pasó?
  \item ¿Has cambiado alguna vez tu forma de escribir al llenar una solicitud,
    postularte a un empleo o hablar con un chatbot o sistema automatizado?
    ¿Qué cambiaste y por qué?
  \item ¿Qué palabras, expresiones o formas de escribir crees que una empresa
    o un sistema automatizado vería ``mal'' en un currículum? Da ejemplos si
    puedes.
  \item Si pudieras decirle algo a quienes diseñan estos sistemas
    automatizados, ¿qué les dirías?
\end{enumerate}

\section*{Cierre}

\begin{quote}
Gracias por participar. Si conoces a otras personas de entre 18 y 30 años que
puedan responder esta encuesta, te agradecemos compartirla
[muestreo por bola de nieve].
\end{quote}

% -----------------------------------------------------------------------------
% Mapa instrumento -> análisis (no se imprime en la versión aplicada;
% conservar aquí como documentación metodológica para el Cap. 3):
%
%   D1-D7   -> caracterización de la muestra (Cap 4, tabla descriptiva)
%   D6-D7   -> variable de cruce socioeconómico para los bloques Likert
%   L1-L4   -> percepción de discriminación lingüística (humana)
%   L5-L9   -> percepción de discriminación tecnológica  (algorítmica)
%   L10-L12 -> exposición y confianza
%   L5, L7  -> ítems invertidos (control de aquiescencia; recodificar)
%   A1-A2   -> corpus: experiencias narradas (temas + sentimiento)
%   A3      -> corpus: léxico estigmatizado percibido -> INSUMO DIRECTO para
%              la tipología de variaciones del benchmark (Fase 3)
%   A4      -> corpus: expectativas normativas (Cap 6, discusión)
% -----------------------------------------------------------------------------

% =============================================================================
%  Apéndice B -- Detalles del corpus controlado
% =============================================================================
\chapter{Detalles del corpus controlado}
\label{ap:corpus}

\section*{Resumen de pares}

[Tabla con estadísticas descriptivas del \emph{benchmark}: número de pares por
tipo de variación, longitud promedio, dispersión.]

\section*{Ejemplos representativos}

[Selección de pares ilustrativos, presentados lado a lado, con la diferencia
de registro explicada y la equivalencia de contenido informativo justificada.]


\backmatter

\printbibliography[heading=bibintoc,title={Referencias}]

\end{document}
